%! AP Statistics — Unit 2, Lesson 2.6 — Exit Ticket
\documentclass[10pt]{article}
\usepackage{apstats-article}
\usepackage{apstats-boxes}

\begin{document}

\pageheader{Unit 2, Lesson 2.6}{Exit Ticket}
\namedateperiod

\vspace{0.18in}

% ── SCENARIO ─────────────────────────────────────────────────────────────────
\begin{tcolorbox}[colback=sky, colframe=navy, boxrule=0.9pt, arc=2mm,
    left=4mm, right=4mm, top=2.5mm, bottom=2.5mm]
\small\textit{A statistics teacher used data from 28 students (who studied between 0 and 8
hours) to produce the regression equation:
$\;\widehat{\text{score}} = 40 + 5(\text{hours studied})$}
\end{tcolorbox}

\vspace{0.16in}

\begin{enumerate}[leftmargin=1.6em, itemsep=0.16in]

  \item Interpret the \textbf{slope} ($5$) in context.\\[8pt]
        \writeline\\[3pt]\writeline

  \item Interpret the \textbf{$y$-intercept} ($40$) in context. Is it meaningful? State why.\\[8pt]
        \writeline\\[3pt]\writeline\\[3pt]\writeline

  \item \textbf{(a)} Predict the score for a student who studied \textbf{6 hours}. Show work
        and classify as interpolation or extrapolation.\\[8pt]
        \writeline\\[3pt]\writeline

        \textbf{(b)} A student claims she studied \textbf{30 hours} and uses the equation to
        predict a score of 190. Why is this prediction not trustworthy?\\[8pt]
        \writeline\\[3pt]\writeline

\end{enumerate}

\vspace{0.14in}

% ── AP-STYLE MC ──────────────────────────────────────────────────────────────
\begin{tcolorbox}[colback=redbg, colframe=redacc, boxrule=0.8pt, arc=2mm,
    left=4mm, right=4mm, top=2.5mm, bottom=2.5mm]
\small\textbf{AP-Style Practice}\enspace
For the regression $\widehat{\text{score}} = 40 + 5(\text{hours})$, which statement correctly
interprets the slope?

\vspace{7pt}
\begin{enumerate}[label=\textbf{(\Alph*)}, leftmargin=2em, itemsep=4pt]
  \item A student who studies 0 hours is predicted to score 40.
  \item The predicted score increases by 40 for each additional hour studied.
  \item For each additional hour studied, the predicted score increases by 5 points.
  \item Students who study more always score exactly 5 points higher.
\end{enumerate}
\end{tcolorbox}

\vspace{0.08in}
\noindent\textcolor{slate}{\small My answer: \blank{1cm} \quad Because: \blank{8.5cm}}

\end{document}
